% =====================================================================
%  Model 2 -- the per-fish latent, as an ordered sequence of operations.
%  Requires: amsmath, booktabs.  (No siunitx, no algorithm packages.)
%  Assumes \dth and \E are defined in the preamble.
% =====================================================================

\subsubsection{The mechanism as a sequence of operations}

\paragraph{Notation.} The mechanism involves four objects, and most of the
difficulty in describing it comes from keeping them apart. They are fixed here
and these names are used consistently below.

\begin{center}
\begin{tabular}{llll}
\toprule
symbol & name & size & what it is \\
\midrule
$z_i$ & \emph{latent} & 2 numbers & who fish $i$ is \\
$Z$ & \emph{latent table} & $15\times 2$ & all fifteen latents, one row per training fish \\
$B_\theta,\,B_\ell,\,B_\tau$ & \emph{loadings} & 2 numbers each & how a latent matters, one per channel \\
$\eta^c_r$ & \emph{offset} & 1 number & what is actually added on bout $r$ \\
$g$ & \emph{gain} & 1 number & multiplies all three loadings at once \\
$\theta$ & network weights & ${\sim}20{,}000$ & shared body and heads, as in Model~1 \\
\bottomrule
\end{tabular}
\end{center}

\noindent The three that matter are read together as a single sentence: the
\emph{latent} says who this fish is, the \emph{loading} says how that matters
for one particular channel, and their dot product is the \emph{offset} --- the
one number actually added to that channel's parameter,
\begin{equation}
\eta^{c}_{r}=B_c^{\top}z_{i(r)},
\qquad c\in\{\theta,\ell,\tau\},
\label{eq:offset}
\end{equation}
where $i(r)$ denotes the fish that bout $r$ belongs to. Two numbers describing
the fish, dotted with two numbers describing the channel, give one number added
to the channel. There are three channels, hence three loadings and three
offsets, but only ever one latent per fish.

The pipeline below has nine stages: 1--4 fit, 5--7 calibrate, 8--9 use.

\paragraph{Stage 1 --- Create the latent table and the loadings.}
Alongside the network weights $\theta$ of Model~1, create the latent table $Z$
with one row of $m=2$ numbers per training fish, and three loadings
$B_\theta,B_\ell,B_\tau\in\mathbb{R}^2$. Thirty-six new numbers in total.
Initialise both from small random draws --- the table at a standard deviation of
$0.5$, somewhat tighter than the prior, and the loadings correspondingly. The
particular values are immaterial with one exception: initialising \emph{both} at
zero is an exact
saddle point, since each one's gradient is proportional to the other, and the
fit would never leave it. The network architecture is unchanged and no layer
receives the fish index as an input.

\paragraph{Stage 2 --- Add the three offsets.}
Take a training bout belonging to fish $i$, with features $x$ and observations
$(\dth,\ell,\tau)$. Compute the shared body $h=\mathrm{ReLU}(Wx+b)$ exactly as
in Model~1; it does not depend on $i$. Look up that fish's latent $z_i$ from the
table, form the three offsets by \eqref{eq:offset}, and add each to one head
output and nothing else:
\[
o^{\text{turn}}_F \mathrel{+}= \eta^{\theta},\qquad
\mu^\ell \mathrel{+}= \eta^{\ell},\qquad
\mu^\tau \mathrel{+}= \eta^{\tau}.
\]
That is the entire forward mechanism. Head variances, mixture shapes and the
stimulus response are untouched, so the latent moves \emph{where} each
conditional sits and never its shape.

\paragraph{The prior is fixed, and never estimated.}
It is worth saying at the outset what is \emph{not} fitted, because in a
classical random-effects model it would be the main object of estimation. The
prior over latents is pinned at $\mathcal{N}(0,I)$ before training begins and is
still exactly that at the end: no mean, no variance and no covariance is ever
updated. There is no such parameter to update --- the penalty in Stage~3 is
$\tfrac12\lVert z\rVert^{2}$ with no variance appearing in it at all.

This is forced by \eqref{eq:offset}, where the latent and the loading enter only
as a product: replacing $z\to cz$ and $B\to B/c$ leaves every offset unchanged,
so a prior variance and a loading length are not separately identifiable and
only their product is. Estimating both is ill-posed, and it fails badly rather
than gracefully; attempted here by expectation-maximisation, the estimated scale
ran from $0.49$ to $27.0$ in five iterations, each re-estimate widening the range
over which the next was read off. We therefore adopt the standard
factor-analytic resolution --- fix the prior and let the loadings carry the
magnitude.

The variance component has not been discarded by this, only renamed. What a
classical formulation would call the between-fish standard deviation is
$\lVert B_c\rVert$ here, which is precisely what \eqref{eq:alg-sd-is-norm}
states. It is estimated in two parts: Stage~4 fits the loadings' directions and
their lengths relative to one another by maximum likelihood, and Stage~7 sets
their common scale by moments, for the reason Stage~6 gives. The prior itself
never moves because it carries no information --- it fixes a scale convention
and nothing more.

\paragraph{Stage 3 --- Form the loss.}
Add to the usual negative log-likelihood a penalty on latents that stray from
the origin, scaled to the per-bout loss. This penalty is the prior
$z\sim\mathcal{N}(0,I)$ written as a cost. Training then proceeds on the same
shuffled-bout minibatches as Model~1, with no change to the optimiser and no
sequence batching.

\paragraph{Stage 4 --- Take one optimiser step.}
The latent table holds ordinary learnable parameters and is updated here and
nowhere else. A single Adam step moves the network weights $\theta$, the latent
table $Z$ and the loadings $B$ simultaneously, roughly $235$ times per epoch.
There is no separate update rule for the latent, no alternating scheme and no
expectation-maximisation. The one distinction is that $Z$ occupies its own
parameter group with weight decay disabled, since the penalty of Stage~3 already
\emph{is} its prior and the optimiser's own decay would apply that prior twice.

The division of labour between latents and loadings follows from which bouts
each gradient sees, and is worth writing out because the summation ranges are
the whole argument. For a minibatch $\mathcal{B}$, with
$\mathcal{B}_c\subseteq\mathcal{B}$ the bouts carrying a valid observation in
channel $c$ and $F$ the number of training fish, the loss of Stage~3 is
\begin{equation}
\mathcal{L}_{\mathcal{B}}
= -\sum_{c}\frac{1}{|\mathcal{B}_c|}\sum_{r\in\mathcal{B}_c}\log p_c(y_r)
\;+\;\frac{1}{2N}\sum_{j=1}^{F}\lVert z_j\rVert^{2},
\end{equation}
where $y_r=(\dth_r,\ell_r,\tau_r)$ is what the fish actually did on bout $r$ and
$p_c$ is that channel's predicted density evaluated at the corresponding
observation. Write
$g^{c}_{r}=\partial\log p_c(y_r)/\partial\eta^{c}_{r}$ for how strongly
bout $r$ wants its channel-$c$ offset raised. Since the latent reaches the
likelihood only through the offset \eqref{eq:offset}, the chain rule gives
\begin{align}
\frac{\partial\mathcal{L}_{\mathcal{B}}}{\partial B_c}
&= -\frac{1}{|\mathcal{B}_c|}\sum_{r\in\mathcal{B}_c} g^{c}_{r}\,z_{i(r)},
\label{eq:grad-B}\\[4pt]
\frac{\partial\mathcal{L}_{\mathcal{B}}}{\partial z_j}
&= -\sum_{c}\frac{1}{|\mathcal{B}_c|}
   \sum_{\substack{r\in\mathcal{B}_c\\ i(r)=j}} g^{c}_{r}\,B_c
\;+\;\frac{z_j}{N}.
\label{eq:grad-z}
\end{align}
The two differ only in what they sum over. Equation~\eqref{eq:grad-z} is
restricted to $i(r)=j$: a fish's latent moves on the evidence of that fish's
bouts alone, so if it turns more than the shared body predicts its latent slides
along $-B_\theta$, independently of every other fish. Equation~\eqref{eq:grad-B}
carries no such restriction and is, up to scale, an empirical covariance between
the residual $g^c$ and the latent $z$ across the whole batch. A loading
therefore lengthens only where the two \emph{co-vary} --- where fish with large
latents systematically want their channel-$c$ offset raised. Had individuals
differed in unrelated ways, their latents would point in unrelated directions,
the products $g^{c}_{r}z_{i(r)}$ would carry no common sign, and the sum would
cancel. The \emph{structure} the loadings carry is in this sense discovered
rather than assigned.

That said, care is needed over what exactly has been discovered, because the
individual latent axes have no meaning and cannot be interpreted. For any
rotation $R$, replacing $z\to Rz$ and $B_c\to RB_c$ leaves every offset
\eqref{eq:offset} unchanged, and leaves the prior unchanged too since
$\mathcal{N}(0,I)$ is rotationally symmetric. The two models are identical in
likelihood, in generation and in every reported metric. Statements of the form
``the first axis is a step-length axis'' are therefore properties of where the
optimiser happened to stop, not of the fitted model, and would not survive a
re-run from a different seed.

What \emph{is} invariant, and so genuinely estimated, is the matrix of inner
products $B_c^{\top}B_{c'}$ --- equivalently the between-fish covariance the
model induces across the three channels. Read as correlations, the fitted model
places the forward tendency and the interval at $-0.79$ (straighter swimmers
bout faster, against $-0.50$ in the real fish) and travel against interval at
$+0.68$ (against $+0.06$ in the real fish), with forward tendency against travel
at $-0.09$, effectively null. Only the first of these is a recovered feature of
the data. The second is a structural artefact: with $m=2$ latent dimensions and
three channels the induced covariance has rank two, so the three correlations
cannot be set independently and fitting the strong turn--interval coupling
forces a spurious travel--interval one. This is the same rank restriction that
leaves the interval spread short in Stage~7, and it is the clearest argument for
a third latent dimension.

The residuals are interpretable, which makes this concrete. Consider the turn
channel, where the offset enters the \emph{logit} of the forward component and
the mixture weights are the softmax $w_k=e^{o_k}/\sum_j e^{o_j}$. Only the
weights depend on $\eta^{\theta}$; the component densities $p_k$ do not. Using
$\partial w_k/\partial\eta^{\theta}=w_k(\delta_{kF}-w_F)$,
\begin{align}
g^{\theta}_{r}
=\frac{\partial\log p_\theta}{\partial\eta^{\theta}_{r}}
&=\frac{1}{p_\theta(y_r)}\sum_k p_k(y_r)\,w_k\bigl(\delta_{kF}-w_F\bigr)
\nonumber\\
&=\frac{w_F\,p_F(y_r)-w_F\,p_\theta(y_r)}{p_\theta(y_r)}
\;=\;r_F(y_r)-w_F ,
\label{eq:resp-minus-weight}
\end{align}
the second line following because the $\delta_{kF}$ term retains only $k=F$
while the $-w_F$ term factors out of $\sum_k w_k p_k = p_\theta$, and the last
step being the definition of the responsibility
$r_F=w_F p_F/p_\theta$, the posterior probability that this observed turn came
from the forward component.

So $g^{\theta}_{r}$ is a prediction error on the mixture weight: the
responsibility the bout turned out to deserve, minus the weight the model had
assigned it beforehand. A near-zero turn is explained only by the forward
component, giving $r_F\approx1$ and a positive residual; a turn of $60^{\circ}$
is explained only by the Gammas, giving $r_F\approx0$ and a negative one; a bout
exactly as ambiguous as predicted gives $r_F=w_F$ and no push at all. This is
not special to this model --- for any softmax-weighted mixture, the derivative of
the log-likelihood with respect to a component's logit is that component's
responsibility minus its weight, which is the same identity underlying the
M-step of expectation-maximisation. Note also that both terms in
\eqref{eq:resp-minus-weight} are probabilities, so $g^{\theta}_{r}$ is bounded
in $(-1,1)$; this is why the turn loading needs no re-centring in Stage~7, where
an unbounded location shift would.

For the log-normal channels, where the offset shifts a location rather than a
logit, the analogous residual is
$g^{\ell}_{r}=\sum_k r_k(\log\ell_r-\mu_k-\eta^{\ell}_{r})/\sigma_k^{2}$, a
responsibility-weighted standardised residual, positive exactly when the fish
travelled further than predicted. So $B_\ell$ lengthens when fish with large
latents consistently overshoot the predicted travel distance, which is the
verbal claim above stated as an estimator.

The prior term in \eqref{eq:grad-z} is the one part that does not depend on the
minibatch: it is evaluated over the whole table at every step, so all $F$
latents are drawn toward the origin whether or not their fish appears in the
batch.

\paragraph{Stage 5 --- Read off the spread the fit produced.}
At generation a latent will be drawn from $\mathcal{N}(0,I)$, so by
\eqref{eq:offset} the offset is a linear map of a standard normal and
\begin{equation}
\mathrm{SD}\bigl(\eta^{c}\bigr)=\lVert B_c\rVert
\label{eq:alg-sd-is-norm}
\end{equation}
exactly. This is the hinge of the whole scheme, and it says something simple:
\emph{the between-fish spread the model will generate in channel $c$ is the
length of that channel's loading}. A longer loading turns the same latent into a
bigger offset, so individuals differ more. The loading's \emph{direction} says
which channels co-vary; its \emph{length} says by how much. For travel distance,
$\lVert B_\ell\rVert=\sqrt{0.171^2+0.101^2}=0.199$, so individuals differ in
median travel by a factor $e^{0.199}=1.22$ per standard deviation.

\paragraph{Stage 6 --- Recognise that the fit does not pin those lengths.}
One would expect the Stage~3 penalty to have controlled the loading lengths.
It has not, and the reason can be stated exactly, without approximation.
Consider rescaling $B_c\to cB_c$ and $z_i\to z_i/c$ for any $c>0$. The latent
reaches the likelihood only through the offset \eqref{eq:offset}, and that
product is unchanged, so the likelihood term is \emph{exactly} invariant. The
penalty, however, becomes
\begin{equation}
\frac{1}{2N}\sum_j\Bigl\lVert \frac{z_j}{c}\Bigr\rVert^{2}
=\frac{1}{c^{2}}\cdot\frac{1}{2N}\sum_j\lVert z_j\rVert^{2},
\label{eq:scale-degeneracy}
\end{equation}
which decreases monotonically in $c$. The loss therefore has a direction along
which the fit is unchanged and the penalty falls without limit: there is no
stationary point along it, gradient descent always has an incentive to push
further, and the penalty can be made arbitrarily small at no cost. Whatever
else the prior does, it cannot constrain the product $B_c^{\top}z_i$, and by
\eqref{eq:alg-sd-is-norm} it therefore cannot constrain the generated spread.

\paragraph{Where the equilibrium sits.} The argument above shows the penalty is
escapable but not where the fit lands, for which the likelihood's shape is
needed. A scalar analogue makes this tractable, and it is obtained from the true
loss in four steps.

First, group the bouts of Stage~4's loss by fish. This is exact --- merely
reordering a sum --- and it is the step that makes everything else possible,
since every bout of a given fish sees the \emph{same} offset:
\begin{equation}
\sum_{r}\log p\bigl(y_r\mid B^{\top}z_{i(r)}\bigr)
=\sum_{i}\ell_i\bigl(B^{\top}z_i\bigr),
\qquad
\ell_i(\eta)\equiv\!\!\sum_{r\in\text{fish }i}\!\!\log p(y_r\mid\eta).
\end{equation}
A whole fish's contribution is thus a function of a single scalar.

Second, $\ell_i$ is a smooth curve in one variable and so attains a maximum
somewhere; write $\delta_i=\arg\max_\eta \ell_i(\eta)$ for its location. This is
a definition, not an assumption: $\delta_i$ is simply a name for the offset that
fish $i$'s own bouts would prefer, and it appears nowhere in the model itself.

Third, expand $\ell_i$ about that maximum. The first derivative vanishes there,
so the quadratic is the leading surviving term,
\begin{equation}
\ell_i(\eta)=\ell_i(\delta_i)-\tfrac{\kappa}{2}\,(\eta-\delta_i)^{2}
+O\!\left((\eta-\delta_i)^{3}\right),
\end{equation}
with $\kappa=-\ell_i''(\delta_i)$ the observed Fisher information for the
offset --- how sharply the peak falls away --- which for $n$ bouts of noise
$\sigma^2$ is $\kappa=n/\sigma^{2}$. This is the standard Laplace
approximation: replace a hill by the parabola through its summit.

Fourth, substitute $\eta=Bz_i$, drop the constants, restore the penalty, and
collapse to one channel and one latent dimension so that $B$ and $z_i$ are
scalars:
\begin{equation}
\mathcal{L}=\frac{\kappa}{2}\sum_i(\delta_i-Bz_i)^2+\frac12\sum_i z_i^2 .
\label{eq:scalar-analogue}
\end{equation}
Each fish is now penalised in proportion to how far its offset sits from its own
preferred value, which is what the true loss says to second order. From
\eqref{eq:scalar-analogue} the two stationarity conditions give
\[
z_i=\frac{a\,\delta_i}{B},\quad a=\frac{\kappa B^2}{\kappa B^2+1}
\qquad\text{and}\qquad
B=\frac{\sum_i z_i\delta_i}{\sum_i z_i^2}=\frac{B}{a},
\]
and the second forces $a=1$: at the joint optimum the prior is inactive and
$Bz_i=\delta_i$, the \emph{unshrunk} per-fish estimate. Since $\delta_i$ is
measured rather than known,
\begin{equation}
\lVert B\rVert^2=\operatorname{Var}(\delta)
=\operatorname{Var}(\delta^{\text{true}})+\operatorname{Var}(\varepsilon)
>\operatorname{Var}(\delta^{\text{true}}),
\end{equation}
the familiar upward bias of a variance component built from unshrunk per-group
estimates. By \eqref{eq:alg-sd-is-norm} the generated population is therefore
too \emph{wide}, and its width must be set by something other than the fit.

Two caveats on how far this reasoning is meant to carry. The quadratic step
drops cubic and higher terms, assumes a common $\kappa$ across fish where the
counts in fact range from $3{,}129$ to $6{,}686$ bouts, and collapses three
channels and two latent dimensions to one apiece; the numerical value of the
shrinkage factor $a$ should be read in that light. The \emph{conclusion},
however, does not rest on any of it --- that the penalty cannot pin the product,
and hence cannot pin the generated spread, follows from the exact invariance
\eqref{eq:scale-degeneracy}. The analogue serves only to locate the equilibrium
that invariance implies.

\paragraph{Stage 7 --- Set the loading lengths, then restore the centre.}
All three loadings are multiplied by one shared number, the gain $g$. Since
\eqref{eq:alg-sd-is-norm} makes each channel's generated spread equal to its
loading's length, multiplying every loading by $g$ multiplies the whole
generated population's spread by $g$. Choosing the gain is therefore choosing
how different from one another the generated fish are.

It is chosen by sweeping a bounded grid and scoring each value on how well the
generated spread matches the real one,
\[
\mathrm{err}(g)=\sum_{j=1}^{3}
\Bigl|\log\bigl(\mathrm{SD}^{\text{gen}}_j / \mathrm{SD}^{\text{real}}_j\bigr)\Bigr| ,
\]
over the across-fish spreads of mean $\log\ell$, mean $\log\tau$ and forward
fraction, computed on training fish only. Log-ratio units are used because the
three statistics live on different scales and an unnormalised sum would let the
smallest dominate.

Note what a single shared gain can and cannot do. It preserves every loading's
direction and the three loadings' lengths \emph{relative to each other};
it changes only their common scale. So it cannot repair a channel whose spread
is wrong relative to the others --- pulling one loading up would drag the other
two with it.

The gain cannot instead be chosen by likelihood. Each training fish sits at
$z_i=\delta_i/B$ by Stage~6, so all of them lie far out along their own loading
direction and a wider prior always scores them better; the marginal likelihood
is monotone in $g$ with no interior optimum, and a leak-free criterion would
require refitting with each fish held out in turn.

Finally, lengthening the loadings would move the pooled marginal as a side
effect, because for a log-normal channel an offset enters the arithmetic mean,
$\E[e^{\eta}]=e^{\lVert B\rVert^2/2}$. Rescaling multiplies every generated
travel and interval by $\exp\bigl(\tfrac12(g^2-1)\lVert B\rVert^2\bigr)$ with
nothing to cancel it, so exactly that amount is subtracted from the
corresponding head biases in the same operation. Changing the \emph{width} of
the population then leaves its \emph{centre} where it was. The turn loading
needs no analogue: its offset tilts a mixture logit rather than a location, so
it is bounded by the softmax, and its pooled drift is held by the gain objective
directly.

\paragraph{Stage 8 --- Score a fish the model has never seen.}
A held-out fish has no row in the latent table, so its latent is integrated out
rather than estimated. Place $Q=49$ candidate latents on a $7\times7$
Gauss--Hermite grid and give each its prior weight. Then, for each bout in turn:
evaluate the bout's log-density under all $Q$ candidates; report as the
predictive density the weighted average over the \emph{current} weights; and only
then multiply each weight by how well its candidate did.

Reporting before updating is what makes this an honest one-step-ahead
prediction: the weights used to predict a bout know only about earlier bouts.
Read plainly, the model does not know which individual it is watching, so it
predicts by averaging over the individuals it still considers possible and
re-weights them after every bout. Early bouts are predicted by a diffuse average
over the population; by the end of a trajectory the weight has concentrated and
predictions come from a nearly identified individual. That progressive
identification is where the likelihood improvement arises, and it is unavailable
to a model without a latent.

Because a latent is constant along a trajectory, each candidate's weight is
simply its prior times the product of its own past likelihoods, so the update is
a cumulative sum in logs rather than a recursion and no weight ever moves
\emph{between} candidates. That is why $49$ candidates cost little. The quantity
reported is the same marginal predictive density any latent-variable model
reports, so Metric~A stays comparable across models.

\paragraph{Stage 9 --- Generate a new individual.}
Draw a latent from $\mathcal{N}(0,I)$ once per trajectory, form the three
offsets by \eqref{eq:offset} and fold them into the head biases before the
rollout starts; since the latent never changes, the per-bout cost is zero. The
draw is taken from the prior and not from the latent table: the generated fish
is a \emph{new} individual from the estimated population rather than a
reconstruction of the held-out one, which is the only choice that does not
condition the generator on the data it is about to be compared against.

\paragraph{What each stage touches.}
Three stages above involve a latent, and only one of them fits anything.

\begin{center}
\begin{tabular}{llll}
\toprule
stage & latent table $Z$ & loadings $B$ & fitted? \\
\midrule
4 (training)   & updated by gradient descent & updated by gradient descent & yes \\
7 (calibration)& untouched                   & all three multiplied by $g$ & by moments \\
8 (scoring)    & never consulted             & read only                   & no \\
9 (generation) & never consulted             & read only                   & no \\
\bottomrule
\end{tabular}
\end{center}

\noindent After Stage~4 the latent table is never read again. Stage~7 changes
only the loadings; Stage~8 averages over latents laid out on a fixed grid; and
Stage~9 draws a fresh one. In particular, at score time the model never
estimates a latent for the held-out fish: it holds $49$ fixed candidates and
revises only \emph{how strongly it believes each one}, which is a posterior over
a grid rather than a fitted value, and which is then integrated back out of the
reported number. Nothing about the evaluation fish is ever written into a
parameter.

% =====================================================================
